% Reviewed duplicate FF.31--FF.32; current statements reconciled with FF.35. Do not input.
\begin{claim}{FF.31}{Proved}\label{clm:FFcyclotomictest}
For every prime \(p\), put
\[
                   U_p=T_p(\mu_{p^\infty})
                       \subset\overline{\mathbb Q}.     \tag{FF.31.1}
\]
This countable, transcendence-degree-zero field has
henselian value group \(\mathbb Q\), residue field
\(\overline{\mathbb F}_p\), and contains every root
of unity. For every finite extension \(E/U_p\),
\[
           C_{\operatorname{Aut}(\overline{\mathbb Q})}
                         (G_E)=1.                       \tag{FF.31.2}
\]
It lies on none of the sufficient field lists
preceding this claim, including the bounded-wild class,
and none of the negative centralizer-ghost lists.
Thus \(U_5\) was a specified next test with the
smallest possible cardinality and transcendence degree
in characteristic zero; no inclusion-minimality is
claimed. Claim~\ref{clm:FFcyclotomicreconstruct}
subsequently proves reconstruction for every finite
extension of \(U_p\).
\end{claim}
\begin{proof}
Henselianity and the residue statement follow by
algebraic extension from \(K_p\). The tame
uniformizer roots supply all prime-to-\(p\)
denominators in the value group, while
\(v(\zeta_{p^n}-1)=1/(p^{n-1}(p-1))\)
supplies arbitrarily large \(p\)-denominators.
Their sums generate \(\mathbb Q\).

We first record the group input precisely.
For a finite \(p\)-adic field \(M\), its wild
inertia \(P_M\) is free pro-\(p\) of countably
infinite rank by
\citet[Theorem~2(ii), p.~463]{FF_Iwasawa1955}.
Closed subgroups of a free pro-\(p\) group are
free pro-\(p\); in a finite-rank free pro-\(p\)
group the centralizer of a nontrivial element
is procyclic. The latter statement and the
subgroup property are explicitly used and
proved/cited in
\citet[Theorem~5.1, proof, Case~I]{FF_KochloukovaZalesskii2011}.
Consequently, if \(H\subset P_M\) is nonabelian,
then \(C_{P_M}(H)=1\). Indeed if a nontrivial
\(c\) commuted with noncommuting
\(h_1,h_2\in H\), the closed subgroup
\(\langle c,h_1,h_2\rangle\) would be a
finite-rank free pro-\(p\) group in which both
\(h_i\) belong to the procyclic centralizer
of \(c\), a contradiction.

For
\(H=\ker(\chi_p:P_M\to\mathbb Z_p^\times)\),
nonabelianity is explicit. Choose four free
basis elements and map them to
\(1+e_{12},1+e_{23},1+e_{34},1+e_{45}\)
in the upper unitriangular group
\(\operatorname{UT}_5(\mathbb F_p)\).
Their first and second pair commutators map
to \(1+e_{13}\) and \(1+e_{35}\), whose
commutator is \(1+e_{15}\ne1\).
Both pair commutators belong to \(H\), since
\(\chi_p\) is abelian. This works also for
\(p=2\).

Now let \(E/U_p\) be finite. Choose a finite
extension \(L_0/K_p\) inside \(E\) such that
\(E=L_0U_p\), and then a number field model
with completion \(M\) for \(L_0\), as in
Claim~\ref{clm:FFstrictreconstruct}. We have
\[
 G_{L_0T_p}=P_M,
 \qquad G_E=\ker(\chi_p|_{P_M})=:H,
 \qquad H\triangleleft G_M.                            \tag{FF.31.3}
\]
In particular \(E\) is not algebraically closed.
Let \(C=C_{G_M}(H)\). Since \(H\) is normal,
\(C\) is normal in \(G_M\); the preceding
paragraph gives \(C\cap P_M=1\).
As \(P_M\) is normal too,
\([C,P_M]\subset C\cap P_M=1\).
But \(L_0T_p/T_p\) is finite, and
Claim~\ref{clm:FFtamereconstruct}, followed by
the necessary centralizer condition of
Claim~\ref{clm:FF6}, gives
\(C_{G_{\mathbb Q}}(P_M)=1\).
Therefore \(C=1\).

Suppose finally that
\(\sigma\in\operatorname{Aut}(\overline{\mathbb Q})
=G_{\mathbb Q}\) centralizes \(G_E\).
It preserves \(E\). Its transported rank-one
henselian valuation must equal the chosen one:
otherwise the two-valuation argument of
Claim~\ref{clm:FFstricttest} would make \(E\)
algebraically closed. Uniqueness of prolongation
from a henselian field then puts \(\sigma\)
in the chosen decomposition group
\(G_{\mathbb Q_p}\subset G_{\mathbb Q}\).
Since \(G_M\) is open, some positive power
\(\sigma^r\) belongs to \(G_M\), hence to
\(C=1\). The group \(G_{\mathbb Q_p}\) has
no torsion: its wild inertia is torsion-free
free pro-\(p\), and its tame quotient is an
extension of \(\widehat{\mathbb Z}\) by
\(\prod_{\ell\ne p}\mathbb Z_\ell\), both
torsion-free. Thus \(\sigma=1\), proving
(FF.31.2).

For completeness the list exclusions do not
rely only on failure of a sufficient hypothesis.
The two-valuation normal-closure proof of
Claim~\ref{clm:FFstricttest} applies to every
finite extension of \(U_p\). Thus its normal
closure over any number field of definition
is \(\overline{\mathbb Q}\), and a distinct
valuation-conjugate has compositum
\(\overline{\mathbb Q}\). Since its absolute
Galois group is infinite by (FF.31.3), the same
open-subgroup argument excludes being finite
over a proper normal algebraic field. It also
cannot be contained in a proper normal algebraic
extension of a number field: first adjoin that
number field to \(U_p\), a finite extension,
and apply the normal-closure assertion.
Positive-transcendence-degree function-field
classes are excluded by transcendence degree
zero. The group \(\mu_\infty\subset U_p^\times\)
is nontrivial divisible, excluding the
Kummer-faithful sub-\(q\)-adic classes. For
any valuation of residue characteristic \(q\),
its elements \(\zeta_{q^n}-1\) have unbounded
\(q\)-denominators, excluding the bounded-wild
classes at every prime, under any field embedding.
Finally (FF.31.2) excludes all negative
centralizer-ghost families. The proof of
reconstruction in Claim~\ref{clm:FFcyclotomicreconstruct}
requires the separate treatment of invisible
functionals developed next.
\end{proof}

\begin{claim}{FF.32}{Proved}\label{clm:FFcyclotomicdetector}
The functional-detection assertion of
Claim~\ref{clm:FFpdetector} fails over the
explicit field \(U_5\). More precisely there
is an elliptic curve \(A/U_5\) and a nonzero
\(G_{U_5}\)-equivariant functional
\[
               V_5A\longrightarrow\mathbb Q_5(1)
\]
whose push-forward of every point Kummer class,
over every finite extension of \(U_5\), is zero.
This is not a counterexample to
\(\mathcal R(U_5)\).
\end{claim}
\begin{proof}
Take
\(A:y^2=x^3-x\), whose discriminant \(64\)
is a unit at \(5\), and choose
\(i\in K_5\) with \(i^2=-1\) and
\(i\equiv2\pmod5\). The automorphism
\([i](x,y)=(-x,iy)\) defines the usual
\(\mathbb Z[i]\)-action. Put
\(\pi=2+i\) and \(\bar\pi=2-i\).
Then \([\pi]\) is an isogeny of degree
\(N(2+i)=5\), and its tangent map is
multiplication by \(2+i\), a unit at \(5\).
Hence \([\pi^n]\) is a finite etale isogeny
of the good elliptic model over
\(\mathcal O_{K_5}\), of degree \(5^n\).
Its kernel is constant over this strictly
henselian ring. The idempotents of
\(\mathbb Z[i]\otimes\mathbb Z_5
\simeq\mathbb Z_5\times\mathbb Z_5\)
split \(T_5A\); the \(\pi\)-primary factor
is the constant etale Tate module
\(T_\pi A\simeq\mathbb Z_5\).
Projection to this factor is nonzero.

For any finite \(L/U_5\), properness extends
\(a\in A(L)\) to the valuation ring
\(\mathcal O_L\). This ring is strictly
henselian, even though it is not discrete:
it is henselian and its residue field remains
\(\overline{\mathbb F}_5\). Pulling back
\([\pi^n]\) along the section gives a finite
etale scheme over \(\mathcal O_L\), hence
it splits. Thus \(a\) has a \(\pi^n\)-division
point in \(A(L)\) for every \(n\). Applying
the same splitting successively to each chosen
division point gives a compatible division
sequence in \(A(L)\). Its full Tate-module
Kummer cocycle is therefore zero, so the
Kummer class projected to \(T_\pi A\) is zero.
Since \(5=\pi\bar\pi\), the other factor
acts invertibly on the \(\pi\)-primary part,
so this is exactly the projection of the
ordinary \(5\)-adic point Kummer class.
But \(U_5\) contains all \(5\)-power roots
of unity. Its cyclotomic representation is
therefore trivial, and any chosen nonzero
identification
\(T_\pi A\otimes\mathbb Q_5\simeq
\mathbb Q_5(1)\) is \(G_{U_5}\)-equivariant.
This is the required undetected functional.

No pinned topos map realizing that functional
has been constructed. In particular a pinned
map with constant closed-point function must
factor through that closed point by the
Jacobson argument of
Claim~\ref{clm:FFstrictdetection}; such a map
cannot realize this nonzero geometric
functional. Thus the detector counterexample
does not decide reconstruction or refute
centralizer sufficiency.
\end{proof}

